% =============================================================================
%  Resumen Capítulo 5 — CLASIFICACIÓN
%  Libro: Evsukoff, A. (2020). Inteligência Computacional: Fundamentos e
%         Aplicações. Río de Janeiro: E-papers.
% =============================================================================
% =============================================================================
%  preambulo.tex — Preámbulo de los resúmenes en español
%  Evsukoff, A. (2020). Inteligência Computacional. Rio de Janeiro: E-papers.
% =============================================================================
\documentclass[11pt,a4paper]{article}

\usepackage[T1]{fontenc}
\usepackage[utf8]{inputenc}
\usepackage[spanish,es-tabla]{babel}
\usepackage{lmodern}
\usepackage[a4paper,margin=2.4cm]{geometry}
\usepackage{amsmath,amssymb,amsthm,mathtools}
\usepackage{bm}
\usepackage{enumitem}
\usepackage{xcolor}
\usepackage{tcolorbox}
\tcbuselibrary{skins,breakable}
\usepackage{hyperref}
\usepackage{microtype}
\usepackage{titlesec}
\usepackage{fancyhdr}
\usepackage{booktabs}
\usepackage{array}

\hypersetup{
  colorlinks=true,
  linkcolor=black!70,
  urlcolor=blue!60!black,
  citecolor=black!70,
  pdfborder={0 0 0}
}

\titleformat{\section}{\Large\bfseries\sffamily\color{black!85}}{\thesection}{0.7em}{}
\titleformat{\subsection}{\large\bfseries\sffamily\color{black!75}}{\thesubsection}{0.6em}{}

\pagestyle{fancy}
\fancyhf{}
\fancyhead[L]{\small\itshape Resumen — Inteligencia Computacional (Evsukoff, 2020)}
\fancyhead[R]{\small\thepage}
\renewcommand{\headrulewidth}{0.3pt}

\newtcolorbox{conceito}[1][]{
  colback=blue!4, colframe=blue!55!black, boxrule=0.6pt,
  arc=2mm, left=6pt, right=6pt, top=4pt, bottom=4pt,
  title={\bfseries Concepto clave}, fonttitle=\sffamily\bfseries,
  breakable, #1
}
\newtcolorbox{formula}[1][]{
  colback=gray!6, colframe=gray!50, boxrule=0.4pt,
  arc=1.5mm, left=6pt, right=6pt, top=3pt, bottom=3pt,
  breakable, #1
}
\newtcolorbox{aplicacao}[1][]{
  colback=green!4, colframe=green!45!black, boxrule=0.5pt,
  arc=2mm, left=6pt, right=6pt, top=3pt, bottom=3pt,
  title={\bfseries Aplicación}, fonttitle=\sffamily\bfseries,
  breakable, #1
}

\newcommand{\R}{\mathbb{R}}
\newcommand{\E}{\mathbb{E}}
\newcommand{\Var}{\operatorname{Var}}
\newcommand{\Cov}{\operatorname{Cov}}
\newcommand{\diag}{\operatorname{diag}}
\newcommand{\argmin}{\operatorname*{arg\,min}}
\newcommand{\argmax}{\operatorname*{arg\,max}}
\newcommand{\T}{^{\!\top}}
\DeclareMathOperator{\sign}{sign}
\DeclareMathOperator{\tr}{tr}


\title{\sffamily\bfseries Capítulo 5 — Clasificación\\
       \large Decisión bayesiana, modelos lineales y evaluación de clasificadores}
\author{Resumen de estudio}
\date{}

\begin{document}
\maketitle
\thispagestyle{fancy}

\begin{conceito}
La \textbf{clasificación supervisada} construye un modelo capaz de
predecir la clase correcta $C_j$ de un registro a partir de su vector
de variables de entrada $\mathbf{x}\in\R^p$. El clasificador particiona
(implícita o explícitamente) el espacio de entrada en regiones
asociadas a las clases; la frontera entre regiones es la
\emph{superficie de decisión}. Los errores aparecen cuando registros
próximos pertenecen a clases distintas (ruido, variables no observadas,
solapamiento intrínseco). El capítulo formaliza la clasificación a
partir de la \textbf{teoría bayesiana de decisión} y muestra cómo
derivar, desde este marco, modelos lineales (mínimos cuadrados,
regresión logística).
\end{conceito}

\section{Clasificación Bayesiana}

\subsection{Decisión bayesiana}

El Teorema de Bayes relaciona la \emph{probabilidad a posteriori} con
la \emph{distribución condicional} y la \emph{probabilidad a priori}:
\begin{formula}
\[
P(C_i\mid\mathbf{x})=\frac{p(\mathbf{x}\mid C_i)\,P(C_i)}{p(\mathbf{x})},
\qquad p(\mathbf{x})=\sum_{i=1}^{m}p(\mathbf{x}\mid C_i)P(C_i).
\]
\end{formula}
Interpretación: el conocimiento del problema (probabilidad condicional)
es \emph{invertido} por la observación $\mathbf{x}$ --- razonamiento
abductivo.

\paragraph{Decisión por mínimo error.}
Se asigna $\mathbf{x}$ a la clase de mayor posteriori. Las regiones
$\mathcal{R}_i=\{\mathbf{x}:P(C_i\mid\mathbf{x})\geq P(C_j\mid\mathbf{x})\;\forall j\}$
\textbf{minimizan} la probabilidad de error
$P(E)=\sum_i\int_{\mathcal{R}_j}p(\mathbf{x}\mid C_i)P(C_i)\,d\mathbf{x}$.
En el punto de frontera:
$p(\mathbf{x}\mid C_1)/p(\mathbf{x}\mid C_2)=P(C_2)/P(C_1)$ --- la
clase más frecuente ``empuja'' la frontera contra la menos frecuente.

\paragraph{Decisión por mínimo riesgo condicional.}
Cuando los errores tienen costos $\lambda_{ij}$ asimétricos (p.\,ej.,
diagnóstico de enfermedad grave: un falso negativo es mucho peor que
un falso positivo),
\[
R(C_i\mid\mathbf{x})=\sum_{j}\lambda_{ij}P(C_j\mid\mathbf{x}),
\]
y se decide por el \emph{menor riesgo}. En problemas desbalanceados,
los costos pueden usarse para compensar la frecuencia de las clases.

\subsection{Clasificador bayesiano simple (naïve Bayes)}

Asume \textbf{independencia condicional} entre las variables dada la
clase:
\[
p(\mathbf{x}(t)\mid C_j)=\prod_{i=1}^{p}p(x_i(t)\mid C_j),
\quad\text{usualmente }p(x_i\mid C_j)\sim N(\hat\mu_{ij},\hat\sigma_{ij}^2).
\]
Se estiman apenas $2p$ parámetros por clase (vs.\ $p+p(p+1)/2$ de la
normal multivariada). Ajuste muy rápido, lidia bien con variables
numéricas y categóricas mezcladas, tolerante a valores faltantes.

\textbf{Limitación:} no captura correlaciones. Cuando las clases son
linealmente separables y las variables están poco correlacionadas, el
desempeño es bueno (Figura~5.5). En problemas \emph{no} linealmente
separables con fuerte correlación intra-clase, el clasificador falla
(Figura~5.6); aplicar ACP previamente puede decorrelacionar y salvar
el modelo (Figura~5.7).

Aplicaciones típicas: clasificación de documentos, filtros de
\emph{spam}, análisis de sentimientos, componente en \emph{ensembles}.

\subsection{Clasificador bayesiano cuadrático}

Adopta $p(\mathbf{x}\mid C_i)=N(\hat{\boldsymbol\mu}_i,\hat{\boldsymbol\Sigma}_i)$
multivariada \textbf{sin} hipótesis de independencia. La función
discriminante es
\[
g_i(\mathbf{x})=-\tfrac12(\mathbf{x}-\hat{\boldsymbol\mu}_i)
\hat{\boldsymbol\Sigma}_i^{-1}(\mathbf{x}-\hat{\boldsymbol\mu}_i)\T
-\tfrac12\ln|\hat{\boldsymbol\Sigma}_i|+\ln P(C_i),
\]
dominada por término cuadrático: \emph{superficies de decisión
cuadráticas} (cónicas). Cuando $\hat{\boldsymbol\Sigma}_i=\hat{\boldsymbol\Sigma}$
para todas las clases, el término cuadrático es común y el
clasificador se vuelve \textbf{lineal}:
\[
g_i(\mathbf{x})=\hat{\mathbf{x}}\,\boldsymbol\theta_i\T,\qquad
\boldsymbol\theta_{i,1}=\hat{\boldsymbol\Sigma}^{-1}\hat{\boldsymbol\mu}_i\T.
\]
En alta dimensión con pocos registros, es prudente reducir
dimensionalidad (ACP/ADL) antes de estimar las covarianzas --- por
ejemplo, en \emph{Wine} (13~variables, 178~registros), el ADL reduce a
2~variables ($m-1=2$) y estabiliza el ajuste.

\section{Modelos Lineales de Clasificación}

\subsection{Mínimos cuadrados (MMC)}

Se define la \textbf{variable indicadora bipolar}
$v_i(t)=+1$ si $y(t)=i$, $-1$ en caso contrario. El modelo lineal
estima $\hat v_i(t)=\hat{\mathbf{x}}(t)\boldsymbol\theta_i\T$
minimizando el ECM:
\[
\boldsymbol\theta_i^{*}=
(\hat{\mathbf{X}}\T\hat{\mathbf{X}})^{-1}\hat{\mathbf{X}}\T\mathbf{V}_i.
\]
Toda la maquinaria del Capítulo~3 (regularización de Tikhonov, PCA,
sesgo--varianza) se aplica directamente. El modelo lineal de grado~1
genera frontera hiperplana; grado~2, frontera cuadrática.

\subsection{Mínimos cuadrados ponderados (MMCP)}

Se atribuye peso $w(t)$ a cada registro:
\begin{formula}
\[
\boldsymbol\theta_i^{*}=
(\hat{\mathbf{X}}\T\mathbf{W}\hat{\mathbf{X}})^{-1}
\hat{\mathbf{X}}\T\mathbf{W}\mathbf{V}_i,\quad
\mathbf{W}=\diag(\mathbf{w}).
\]
\end{formula}
Para neutralizar el desbalanceo, se sugiere $w(t)=1-N_i/N$ para
$\mathbf{x}(t)\in C_i$ (peso inverso a la frecuencia). Equivale, en la
visión bayesiana, a remover la influencia de la \emph{priori}
$P(C_i)$. Resultado: error global mayor, pero mejor recuperación de la
clase minoritaria.

\subsection{Regresión logística}

La función discriminante es la \textbf{sigmoide} aplicada al modelo
lineal:
\[
g(\mathbf{x},\boldsymbol\theta)=
\frac{1}{1+\exp(-\hat{\mathbf{x}}\,\boldsymbol\theta\T)}\;\approx\; P(C\mid\mathbf{x}),
\]
lo que devuelve una \emph{interpretación probabilística}. Modelando
$P(v=1\mid\mathbf{x})$ como Bernoulli, la verosimilitud es
\[
L(\boldsymbol\theta)=\prod_{t=1}^{N}\hat v(t)^{v(t)}(1-\hat v(t))^{1-v(t)},
\]
y el ajuste minimiza la \textbf{log-verosimilitud negativa} (entropía
cruzada binaria):
\[
\ell(\boldsymbol\theta)=-\sum_{t=1}^{N}\bigl[v(t)\log\hat v(t)+(1-v(t))\log(1-\hat v(t))\bigr].
\]
Como $\boldsymbol\theta$ aparece no linealmente, se resuelve por
Newton--Raphson (algoritmo IRLS, \emph{iteratively reweighted least
squares}) o por gradiente. La extensión multiclase es la
\textbf{regresión softmax} (Capítulo~9).

\section{Evaluación de Clasificadores}

\subsection{Estadísticas a partir de la matriz de confusión}

\begin{center}
\begin{tabular}{c|cc|c}
\toprule
 & $\hat C_1$ & $\hat C_2$ & total \\
\midrule
$C_1$ & VP & FN & $N_1$ \\
$C_2$ & FP & VN & $N_2$ \\
\midrule
total & $\hat N_1$ & $\hat N_2$ & $N$ \\
\bottomrule
\end{tabular}
\end{center}

\begin{itemize}[leftmargin=1.2cm]
  \item \textbf{Exactitud:} $\mathrm{ACC}=(\mathrm{VP}+\mathrm{VN})/N$.
  \item \textbf{Error:} $\mathrm{ERR}=1-\mathrm{ACC}$.
  \item \textbf{Precisión}
        $\mathrm{PRE}(C_1)=\mathrm{VP}/(\mathrm{VP}+\mathrm{FP})$:
        confiabilidad de una predicción positiva.
  \item \textbf{\emph{Recall}/sensibilidad}
        $\mathrm{REC}(C_1)=\mathrm{VP}/(\mathrm{VP}+\mathrm{FN})$:
        proporción de positivos reales capturados.
  \item \textbf{Medida F:}
        $F(C_i)=2\,\mathrm{PRE}\cdot\mathrm{REC}/(\mathrm{PRE}+\mathrm{REC})$
        --- media armónica.
  \item \textbf{Estadística kappa}
        $\kappa=(\mathrm{ACC}-P(E))/(1-P(E))$, con
        $P(E)=P(C_1)P(\hat C_1)+P(C_2)P(\hat C_2)$. Mide acuerdo por
        encima del esperado por azar ($\kappa\geq 0{,}8$ excelente,
        $0{,}6$--$0{,}8$ razonable, $\leq 0$ peor que aleatorio).
\end{itemize}

\subsection{Espacio ROC y AUC}

Ejes: $\mathrm{TFP}=\mathrm{FP}/(\mathrm{FP}+\mathrm{VN})$ y
$\mathrm{TVP}=\mathrm{REC}(C_1)$. El punto $(0,1)$ es el clasificador
ideal; la diagonal corresponde a $\kappa=0$. El \textbf{área bajo la
curva ROC} (resumida en un único punto) es
\[
\mathrm{AUC}=\frac{\mathrm{REC}(C_1)+\mathrm{REC}(C_2)}{2}.
\]
Para múltiples clases, ponderar por las \emph{prioris}:
$\mathrm{AUC}=\sum_i \mathrm{AUC}(C_i)P(C_i)$. La AUC es la métrica
recomendada en problemas \emph{desbalanceados}, pues la exactitud
tiende a favorecer la clase mayoritaria.

\section{Aplicaciones}

\begin{aplicacao}
\textbf{Benchmarks UCI} (10 conjuntos: Iris, Balance, Diabetes Pima,
Cancer, Statlog Heart, Wine, Parkinson, Ionosphere, Spambase, Sonar).
Comparados Bayes Naïve, Bayes Cuadrático, regresión logística,
modelo lineal de 1\textsuperscript{er} y 2\textsuperscript{do} orden
en validación cruzada de 10 ciclos. El \textbf{modelo lineal de
2\textsuperscript{do} orden} (M.\,Lin.\,2) obtiene el mejor
\emph{ranking} por el test no paramétrico de Friedman \emph{aligned
rank} ($p\approx 0{,}08$ tanto en ACC como en AUC), confirmando que el
MMC es más robusto que la estimación de parámetros del clasificador
bayesiano en conjuntos pequeños.
\end{aplicacao}

\begin{aplicacao}
\textbf{Análisis de \emph{churn} en telefonía} (Cister, 2007). Base
con 1\,141 clientes corporativos de alto valor, 6~meses de consumo
medio ($x_1,\dots,x_6$). Sólo el 7{,}17\% canceló (clase $C_2$,
\emph{minoritaria y más relevante}). \textbf{Sin ponderación}, el
modelo bayesiano simple alcanza ACC$\,=92{,}82\%$ pero AUC$\,=0{,}51$
--- prácticamente aleatorio, pues acierta todo en la clase
mayoritaria. \textbf{Con ponderación} por el complemento de la
\emph{priori}, ACC baja pero AUC sube a $0{,}79$--$0{,}89$,
recuperando la clase de cancelaciones. Lección: en problemas altamente
desbalanceados, ACC \emph{engaña}; usar AUC y \emph{recall} de la
clase minoritaria.
\end{aplicacao}

\section{Comentarios finales}

\begin{enumerate}[leftmargin=1.2cm]
  \item La regla de Bayes es el esqueleto conceptual de la
        clasificación: aunque el clasificador final no sea
        explícitamente bayesiano, los efectos de la \emph{priori} y de
        la distribución condicional impregnan cualquier modelo ajustado
        por minimización del error.
  \item En general, \textbf{estimar la función discriminante} (modelos
        lineales, logística, SVM, redes neuronales) es más fácil que
        estimar las distribuciones condicionales completas en alta
        dimensión.
  \item La \textbf{validación cruzada} y la elección cuidadosa de
        métricas (AUC, $\kappa$, F) son indispensables en problemas
        desbalanceados.
  \item Tópicos correlativos no cubiertos: árboles de decisión
        (Quinlan), redes bayesianas (Koller \& Friedman), SVM (Vapnik)
        y redes neuronales profundas --- éstas tratadas en los
        Capítulos 7, 8 y 9.
  \item Para el desbalance, además de ponderar registros, hay técnicas
        de \textbf{re-muestreo} como el algoritmo SMOTE, muy eficaz
        en conjunto con árboles de decisión.
\end{enumerate}

\end{document}
